\documentclass[12pt]{article}
\usepackage{amsmath}
\usepackage{geometry}
\geometry{letterpaper, margin=1in}
\usepackage{fancyhdr}
\usepackage{setspace}
\usepackage{hyperref}

\pagestyle{fancy}
\fancyhf{}
\rhead{33-120 Science and Science Fiction}
\lhead{SETI Team Project (25 points)}
\setlength{\headheight}{15pt}

\begin{document}

\vspace{1cm}
{\large{\noindent Imagining a Realistic Scenario for First Contact}}

\noindent\rule{6.5in}{0.4pt}

\noindent Now that we have covered the search for extraterrestrial intelligence
in class, your group is equipped to imagine a realistic scenario for \emph{first
contact} with extraterrestrial life. While we don't know exactly how first
contact would play out since will are pretty sure that it hasn't happened
(..yet) your group should work to maintain scientific plausibility throughout
your scenario. 

\vspace{0.3cm}
\noindent As a starting point, you should assume that the contact is \emph{not
an invasion}. However, it is reasonable to assume that any alien intentions are
unclear until the scenario plays out. Throughout the project, you should
consider realistic factors like communication barriers (both physical and
social) as well as social response (both general public as well as
administrative/governmental). 

\vspace{0.3cm}
\noindent Your group should develop the scenario together in detail and discuss
how you will respond to each of the prompts below. Each team member should be
assigned a contribution and your responses to each prompt should include a
description of who contributed to that prompt. 

\vspace{0.3cm}
\noindent Your responses should be gathered into a single document (Google Docs
is a good starting point for this!) for submission to Canvas. Only one
submission per group is necessary. 

\vspace{1cm}



{\bf{\noindent Describing Your \emph{First Contact} Scenario}}
\vspace{5pt}

\noindent Describe your first contact scenario in 3-5 paragraphs. Your
description should address the following:
\begin{itemize}
        \item How does first contact occur? (e.g., a distant radio signal, a
        robotic probe, or a physical arrival)
        \item Where does it take place? (Earth, orbit, deep space, somewhere
        else?)
        \item Who is involved first? (Scientists, military, governments, media,
        the public, someone else?)
        \item How does society respond? (Do governments decide to reveal or
        conceal the information? What actions do they take? What about other
        institutions like universities or corporations?)
        \item What challenges arise in communication and understanding (both
        with the aliens as well as within our broader society)?
\end{itemize}


\vspace{0.5cm}
{\bf{\noindent A Detailed Window Into Society's Response}}
\vspace{5pt}
\noindent Choose one of the following formats to imagine how your scenario might
play out in more detail:
    \begin{itemize}
        \item \textbf{Government Briefing Memo}: Write a classified memo
        advising world leaders on how to respond.
        \item \textbf{News Report}: Create a 300-word breaking news article
        covering the first 24 hours of contact.
        \item \textbf{Dialogue Scene}: Write a 1-2 page script of a government
        official speaking with an alien entity.
        \item \textbf{Social Media Thread}: Simulate public reaction with 10
        social media posts.
    \end{itemize}

\vspace{0.5cm}

{\bf{\noindent Reflection}}
\vspace{5pt}
\noindent Address the following questions in a paragraph for each point to
reflect on your scenario:
    \begin{itemize}
        \item What considerations did you take to make your scenario as
        realistic as possible?
        \item What are the largest challenges that would arise in your first
        contact scenario? How does this differ from other scenarios (e.g. if
        yours is on Earth, how might different challenges arise in a different
        setting?)?
        \item How might governments and the public react differently to your scenario?
    \end{itemize}


\end{document}
